\chapter{Metodología}
\label{cap:metodologia}

Este capítulo describe cómo se ha organizado y ejecutado el desarrollo de Linceus Assistant. Se cubren cinco aspectos: la metodología de desarrollo ágil empleada y su adaptación al contexto del proyecto, la gestión del código fuente, los estándares de codificación aplicados, la gestión del trabajo y las tareas, y el esquema de versionado del sistema.

% =====================================================
\section{Metodología de desarrollo: Scrum adaptado}
\label{sec:scrum}
% =====================================================

\subsection{Scrum: fundamentos}

Scrum es un marco de trabajo ágil para el desarrollo iterativo e incremental de productos complejos~\cite{schwaber2020scrum}, encuadrado en los principios establecidos por el Manifiesto Ágil~\cite{beck2001agile}. Su unidad de trabajo es el \textbf{sprint}, una iteración de duración fija (típicamente de una a cuatro semanas) al final de la cual se obtiene un incremento del producto potencialmente entregable. Los roles clásicos son el \textit{Product Owner}, responsable de priorizar el \textit{product backlog}; el \textit{Scrum Master}, que facilita el proceso; y el \textit{Development Team}, que ejecuta el trabajo. Las ceremonias principales son la \textit{Sprint Planning} (planificación del sprint), la \textit{Daily Scrum} (sincronización diaria), la \textit{Sprint Review} (revisión del incremento) y la \textit{Sprint Retrospective} (mejora del proceso).

Una de las características definitorias de Scrum es su flexibilidad: la guía oficial no prescribe prácticas de ingeniería específicas, sino un conjunto mínimo de reglas que el equipo adapta a su contexto. Esto lo hace especialmente adecuado para proyectos de alcance variable como un TFG, donde los requisitos evolucionan con el conocimiento adquirido durante el desarrollo.

\subsection{Adaptación al proyecto individual}

Linceus Assistant es un proyecto desarrollado por una sola persona, lo que hace innecesarias varias ceremonias pensadas para coordinar equipos. A continuación se describen las adaptaciones realizadas respecto al Scrum estándar:

\begin{itemize}
    \item \textbf{Sin Daily Scrum.} La sincronización diaria carece de sentido en un equipo unipersonal. El seguimiento del progreso se realiza de forma continua a través del tablero de tareas y el registro de decisiones técnicas en \texttt{docs/sprints/}.

    \item \textbf{Retrospectivas bajo demanda.} No se celebra una retrospectiva formal al cierre de cada sprint. En su lugar, se realiza una revisión del proceso únicamente cuando se detecta un problema recurrente o una desviación significativa respecto a la planificación. Este enfoque evita la sobrecarga documental sin perder la capacidad de mejora continua.

    \item \textbf{Sprint Planning al final de la Review.} La planificación del sprint siguiente se integra en la sesión de revisión del sprint anterior. Una vez evaluado el incremento entregado y ajustado el backlog, se definen directamente los objetivos del próximo sprint. El flujo real es, por tanto: \textit{Review del sprint} $\rightarrow$ \textit{Actualización del backlog} $\rightarrow$ \textit{Definición de objetivos} $\rightarrow$ \textit{Descomposición en tareas}.

    \item \textbf{Rol único con supervisión académica.} Al no haber equipo, los tres roles de Scrum (Product Owner, Scrum Master, Developer) los asume la misma persona. Las funciones propias del Product Owner (definir prioridades, validar el incremento y orientar el alcance del producto) se ejercen en colaboración con el tutor académico, cuyas aportaciones quedan recogidas en la Bitácora de Reuniones descrita en la Sección~\ref{sec:gestion-trabajo}.
\end{itemize}

\subsection{Sprints, número y duración}

El proyecto se ha organizado en \textbf{8 sprints} a lo largo de aproximadamente diez meses (julio de 2025 a mayo de 2026), con una duración nominal de \textbf{3 a 4 semanas} por sprint. La duración real de cada iteración se ajustó en función de la carga lectiva del cuatrimestre y de la complejidad de la épica en curso, lo que dio lugar a sprints más cortos durante los periodos de exámenes y a un sprint largo (S7, aproximadamente seis semanas) en el que se concentraron las épicas finales, el panel de administración, el despliegue y el piloto con usuarios. La cadencia y la duración aproximada de cada sprint, junto con los objetivos asociados y las horas registradas, se recogen en la Tabla~\ref{tab:sprints} del Capítulo~\ref{cap:planificacion}.

Cada sprint cierra con un incremento funcional verificable y un breve registro de decisiones técnicas (fichero \texttt{docs/sprints/sprint-N.md}) que documenta los cambios de diseño no triviales realizados durante la iteración. La correspondencia entre sprints y épicas del backlog se detalla en el Capítulo~\ref{cap:planificacion}.

% =====================================================
\section{Gestión del código fuente}
\label{sec:git}
% =====================================================

El control de versiones del proyecto se gestiona con \textbf{Git}~\cite{chacon2014pro}, y el repositorio remoto se aloja en \textbf{GitHub}. Esta combinación proporciona trazabilidad completa de los cambios, facilita la revisión del historial de decisiones técnicas y permite integrar el pipeline de despliegue continuo~\cite{humble2010continuous} descrito en la Sección~\ref{subsubsec:despliegue}.

\subsection{Política de ramas}

El repositorio mantiene en todo momento exactamente dos ramas activas: \texttt{main} y una rama de trabajo efímera. Esta restricción deliberada simplifica la gestión del historial y evita la acumulación de ramas obsoletas.

\begin{itemize}
    \item \textbf{\texttt{main}}: rama permanente y de producción. Todo merge a \texttt{main} dispara automáticamente el pipeline de GitHub Actions que despliega la nueva versión en el servidor de DigitalOcean.

    \item \textbf{Rama de trabajo}: efímera, con nombre descriptivo del tipo \texttt{feature/\allowbreak<épica>} o \texttt{fix/\allowbreak<ámbito>} durante el desarrollo activo de una épica (por ejemplo, \texttt{feature/\allowbreak rag-pipeline} o \texttt{fix/\allowbreak horarios-slot}). Al finalizar la épica se crea una rama \texttt{test} puntual para ejecutar la suite de validación completa; una vez superadas las pruebas, se hace merge a \texttt{main} y la rama se elimina. En ningún momento coexisten más de dos ramas.
\end{itemize}

El ciclo completo de una épica es, por tanto: \texttt{main} → \texttt{feature/X} (desarrollo) → \texttt{test} (validación) → \texttt{main} (merge y despliegue).

\subsection{Convenciones de commits}

Los mensajes de commit siguen el estándar \textbf{Conventional Commits}~\cite{conventionalcommits}, con el formato \texttt{<tipo>: <descripción>}. Los tipos empleados con mayor frecuencia son \texttt{feat} (nueva funcionalidad), \texttt{fix} (corrección de errores), \texttt{refactor} (reestructuración sin cambio de comportamiento), \texttt{test} (adición o modificación de pruebas) y \texttt{docs} (documentación). Los commits son frecuentes y de alcance reducido: cada commit representa un cambio coherente y compilable, lo que facilita la bisección del historial mediante \texttt{git bisect}~\cite{chacon2014pro} cuando aparecen regresiones.

% =====================================================
\section{Estándares de codificación}
\label{sec:estandares}
% =====================================================

El código Python del proyecto se ajusta a la guía de estilo \textbf{PEP~8}~\cite{pep8}, verificada automáticamente mediante dos herramientas complementarias: \textbf{pylint}~\cite{pylint} y \textbf{flake8}~\cite{flake8}. Pylint cubre el análisis semántico (detección de errores reales, variables no definidas, llamadas incorrectas) mientras que flake8 se centra en la conformidad de estilo: espacios sobrantes, imports no utilizados, longitud de línea e indentación. Ambas herramientas se configuran con un límite de 120 caracteres por línea, valor que equilibra legibilidad y practicidad para un proyecto con prompts al LLM y consultas SQL que no admiten partición arbitraria.

En los casos en que una desviación de estilo es intencionada (un import que actúa como re-exportación pública, una línea larga que contiene una cadena de texto no partible, o un par de instrucciones que forman una unidad semántica indivisible) se documenta explícitamente con una directiva de supresión en la propia línea. Esto hace que las excepciones sean visibles y razonadas, en lugar de silenciosas.

Para el código JavaScript del widget de chat (\texttt{frontend/}), se aplica \textbf{ESLint 8}~\cite{eslint} con la configuración \texttt{eslint:recommended}, extendida con reglas de estilo básicas (comillas simples, punto y coma obligatorio). La configuración se almacena en \texttt{.eslintrc.json} en la raíz del repositorio y se ejecuta con \texttt{npm run lint}, que lanza \texttt{eslint frontend/} sobre los tres ficheros del directorio (\texttt{chatbot-widget.js}, \texttt{admin.js}, \texttt{login.js}).

% =====================================================
\section{Política de delegación al modelo de lenguaje}
\label{sec:politica-llm}
% =====================================================

A lo largo del proyecto se ha tomado de forma sistemática una decisión transversal que, sin estar atada a ninguna épica concreta, condiciona la implementación de prácticamente todas: \textbf{cuándo resolver una subtarea con una heurística determinista (regex, diccionario, fuzzy matching) y cuándo delegar en el modelo de lenguaje}. Esta sección explicita el criterio aplicado, porque condiciona también el comportamiento, el coste y la fiabilidad del sistema en producción.

\subsection{Criterio: heurísticas baratas primero, LLM cuando aporta}

El criterio operativo es jerárquico. Para cada subtarea (clasificar una pregunta, extraer una entidad, normalizar un nombre, decidir una rama de pipeline) se aplica el siguiente orden:

\begin{enumerate}
    \item \textbf{Diccionarios y regex} cuando el espacio de entrada es cerrado o casi cerrado (lista finita de alias, conjunto pequeño de palabras clave, patrones léxicos identificables como \textit{``grupo N''} o \textit{``coordinador''}).
    \item \textbf{Fuzzy matching}~\cite{cohen2003comparison} (\texttt{rapidfuzz}, \texttt{SequenceMatcher}) cuando el espacio de entrada es cerrado pero admite errores tipográficos o variantes morfológicas (nombres de profesor en BD, tokens de tutorías con typos).
    \item \textbf{LLM} solo cuando el espacio de entrada es abierto (generar SQL a partir de lenguaje natural, redactar respuesta natural a partir de \textit{chunks}) o cuando una heurística previa no ha sido concluyente y el coste del LLM se justifica.
\end{enumerate}

Esta jerarquía no es una preferencia estética: responde a tres restricciones concretas del proyecto. Primero, el \textbf{coste por turno}: cada llamada al modelo cuenta hacia la cuota mensual de Gemini, y multiplicarla por dos o tres llamadas en cada turno conversacional convierte un sistema económico en uno difícil de mantener. Segundo, la \textbf{latencia percibida}: una llamada al modelo añade entre 800~ms y 2~s al turno, y el usuario percibe esa latencia incluso cuando la pregunta era trivial. Tercero, y probablemente más importante, las \textbf{alucinaciones}~\cite{ji2023survey}: cada vez que se delega una decisión binaria al LLM hay una probabilidad no nula de que invente un valor o lo interprete fuera del rango esperado, lo que en un asistente académico puede traducirse en información incorrecta entregada con tono autoritario.

\subsection{Aplicación en el sistema}

El criterio se materializa en patrones recurrentes a lo largo del código. La Tabla~\ref{tab:politica-llm} resume las cuatro variantes principales y dónde aparecen. En todas ellas, la elección entre heurística y LLM no fue casual: respondió a la naturaleza concreta de la subtarea y al balance entre los tres factores anteriores.

\begin{table}[hbtp]
\centering
\small
\begin{tabular}{|p{3.4cm}|>{\raggedright\arraybackslash}p{4.5cm}|>{\raggedright\arraybackslash}p{6cm}|}
\hline
\textbf{Patrón} & \textbf{Subtarea típica} & \textbf{Ejemplo en el proyecto} \\
\hline
Solo heurística & Detección de patrones léxicos cerrados & \texttt{\_detectar\_grupo()} (regex \texttt{grupo N / gN}); \texttt{\_pregunta\_menciona\_rol\_rag()} (lista de seis palabras clave); \texttt{\_RE\_CONTINUACION\_TODOS\_GRUPOS} (regex de continuación elíptica) \\
\hline
Heurística + fuzzy & Resolución de un nombre contra un catálogo finito & Cascada de \texttt{resolver\_asignatura()} (alias $\rightarrow$ regex $\rightarrow$ fuzzy BD); filtrado por similitud en profesores (\texttt{nombre\_normalizado}, umbrales 0,30 y 0,60) \\
\hline
Heurística + LLM como respaldo & Clasificación binaria con casos triviales y ambiguos & Clasificador RAG en \texttt{ActionConsultaEspecifica}: primero palabras clave (\textit{evalúa}, \textit{temario}); si no concluye, LLM a temperatura cero \\
\hline
LLM directo & Espacio de entrada o salida abierto & \textit{Text-to-SQL} (generación de SELECT); redacción natural de la respuesta final; extracción del nombre de profesor cuando NLU falla y hay asignatura como contexto \\
\hline
\end{tabular}
\caption{Patrones de delegación al LLM aplicados en Linceus Assistant.}
\label{tab:politica-llm}
\end{table}

Un ejemplo ilustrativo del segundo patrón es la normalización de la titulación en \texttt{ActionCambiarContexto} (Sección~\ref{sec:impl-contexto}): coincidencia exacta sobre un diccionario, \textit{guard} de tokens discriminantes y, solo si nada de lo anterior funciona, fuzzy matching con umbral conservador. La opción de delegar la normalización al LLM se descartó deliberadamente: un fallo en la titulación contamina toda la sesión inyectando un \texttt{titulacion\_id} incorrecto en consultas posteriores, y el coste de una alucinación silenciosa supera con creces el coste de un fallo de heurística (que el usuario percibe como \textit{``no te he entendido''} y reformula).

\subsection{Limitaciones conocidas}

El precio de esta política es la \textbf{rigidez frente a fraseos atípicos y la dependencia del idioma}. Las heurísticas léxicas (palabras clave, regex de continuación, pronombres demostrativos para invertir prioridades de seguimiento) están escritas en español y no cubren formulaciones que se aparten del registro común del estudiantado. Si el sistema se internacionalizara o se quisiera atender a un público con vocabulario más diverso, varias de estas heurísticas tendrían que migrar a alternativas más flexibles: diccionarios por idioma, similitud semántica con \textit{embeddings} multilingües, o clasificación por LLM a coste de la latencia y de las alucinaciones que actualmente se evitan. Esta tensión se aborda en el Capítulo~\ref{cap:conclusiones} como línea de trabajo futura.

% =====================================================
\section{Gestión del trabajo}
\label{sec:gestion-trabajo}
% =====================================================

El seguimiento del proyecto se articula en torno a la \textbf{Bitácora de Reuniones}, un documento estructurado que recoge de forma acumulativa el desarrollo de cada reunión con el tutor y los objetivos acordados para el sprint siguiente. Tras cada reunión, la bitácora se actualiza a partir de la transcripción de la sesión, extrayendo los puntos más relevantes (decisiones tomadas, problemas detectados, cambios de prioridad) y registrando las tareas pendientes para el siguiente ciclo. Este documento actúa de facto como el \textit{product backlog} del proyecto: la lista de tareas pendientes refleja en todo momento el estado acordado con el tutor, sin necesidad de herramientas adicionales de gestión.

Las reuniones de seguimiento se celebraron por \textbf{Microsoft Teams}, con una duración aproximada de 30~minutos cada una, al inicio de cada sprint. En total se documentan \textbf{8 reuniones de sprint} en el Apéndice~\ref{ap:actas}, además de una reunión inicial de toma de contacto con el tutor (anterior a la adjudicación oficial del TFG) en la que se acordó el ámbito general del trabajo. Esta cadencia, ligada uno a uno con los sprints, garantiza que la priorización del backlog y la dirección del proyecto se revisen de forma sistemática.

Complementariamente, las decisiones técnicas no triviales tomadas durante el desarrollo se documentan en los registros de decisiones de sprint (\texttt{docs/\allowbreak sprints/\allowbreak SN/\allowbreak Registro\_decisiones.md}), que mantienen trazabilidad entre el código y el razonamiento que lo originó.

El tiempo dedicado al proyecto se registra con \textbf{Clockify}, donde cada entrada recibe el nombre \texttt{SX} seguido de un guion y la descripción de la tarea (siendo \texttt{X} el número de sprint), lo que permite agrupar el esfuerzo por sprint y correlacionarlo con el backlog sin necesidad de herramientas adicionales.

% =====================================================
\section{Versionado del sistema}
\label{sec:versionado}
% =====================================================

El sistema sigue un esquema de versionado semántico de tres niveles: \textbf{\texttt{vMAJOR.MINOR.PATCH}}, adaptado al ciclo de vida de un chatbot con modelo entrenado.

\begin{itemize}
    \item \textbf{\texttt{MAJOR}}: se incrementa en dos situaciones: al incorporar una \textbf{nueva épica funcional} (un nuevo dominio de capacidades), o al introducir un \textbf{cambio arquitectónico de envergadura dentro de una épica existente} que afecta a su lógica de forma global (por ejemplo, la migración completa a Text-to-SQL o la introducción del pipeline RAG). Ejemplos reales del proyecto: \texttt{v1.x.x} cubre asignaturas con NLU puro, \texttt{v2.x.x} introduce Text-to-SQL y RAG, \texttt{v3.x.x} añade horarios, \texttt{v4.x.x} añade profesorado, \texttt{v5.x.x} incorpora cambios transversales de infraestructura.

    \item \textbf{\texttt{MINOR}}: se incrementa al añadir \textbf{nueva funcionalidad dentro de una épica existente}, como nuevos intents, nuevos filtros de búsqueda o mejoras significativas en la lógica de recuperación, sin alterar la arquitectura general.

    \item \textbf{\texttt{PATCH}}: se incrementa en correcciones de errores, ajustes de prompts, mejoras de ejemplos NLU o pequeñas optimizaciones que no alteran el comportamiento funcional esperado.
\end{itemize}

El versionado está ligado al modelo Rasa: cada reentrenamiento que modifica los ficheros NLU, stories o reglas genera un nuevo artefacto de modelo (\texttt{models/linceus\_vX.Y.Z.tar.gz}) identificado con la versión del sistema. Esto garantiza que cada versión desplegada en producción es reproducible: el par \{tag de Git, artefacto de modelo\} define unívocamente el estado del sistema en ese punto. El historial completo de versiones se mantiene en \texttt{docs/tabla\_versiones.md}.
